\documentclass[10pt,aspectratio=169]{beamer}
\usepackage{poly}
\usepackage[useregional]{datetime2}
\usepackage{graphicx}  % for images
\usepackage{xcolor}    % for color
\usepackage{minted}    % for better code highlighting
\usepackage[T1]{fontenc}  % for better font encoding
\usepackage{inconsolata}  % fallback monospace font

% Minted configuration
\setminted{
    frame=none,
    fontsize=\small,
    linenos=true,
    numbersep=12pt,
    bgcolor=white,
    style=friendly,
    autogobble,
    breaklines,
    baselinestretch=1.2
}

% Add slide numbers to the footer
\setbeamertemplate{footline}[frame number]
\setbeamertemplate{navigation symbols}{}  % Remove navigation symbols

% Define a centered minted environment
\newenvironment{centeredcode}{\begin{center}}{\end{center}}

\title{Operating System (IF2223)}
\subtitle{Memory Management 2 (Paging)}
\author{Martin C.T. Manullang, Eko Dwi Nugroho, Ilham Firman Ashari}
\institute[COMP]{Teknik Informatika, ITERA}
\date{2025}

\begin{document}
\maketitle

\section{Konsep Dasar Paging}

\begin{frame}{Pengertian Paging}
    \framesubtitle{Teknik Manajemen Memori}
    \begin{columns}[T]
        \column{0.48\textwidth}
            \begin{block}{Definisi}
                Paging adalah teknik manajemen memori yang memungkinkan alamat fisik memori tidak harus bersambungan (noncontiguous) saat mengalokasikan ruang untuk proses.
            \end{block}
            
            \begin{itemize}
                \item Mengatasi fragmentasi eksternal
                \item Memungkinkan penggunaan memori fisik lebih efisien
                \item Mendukung implementasi memori virtual
            \end{itemize}
            
        \column{0.48\textwidth}
            \begin{alertblock}{Kelebihan Paging}
                \begin{itemize}
                    \item Tidak ada fragmentasi eksternal
                    \item Mengurangi overhead pemeliharaan memori
                    \item Menyederhanakan alokasi memori
                    \item Mendukung sharing kode antar proses
                \end{itemize}
            \end{alertblock}
    \end{columns}
\end{frame}

\begin{frame}{Page dan Frame}
    \framesubtitle{Komponen Utama Sistem Paging}
    
    \begin{columns}[T]
        \column{0.48\textwidth}
            \begin{block}{Page}
                Unit dari memori logis (ruang alamat proses) dengan ukuran tetap.
                \begin{itemize}
                    \item Biasanya berukuran 4KB - 16KB
                    \item Diidentifikasi dengan page number
                    \item Seluruh page harus termuat dalam frame
                \end{itemize}
            \end{block}
            
        \column{0.48\textwidth}
            \begin{block}{Frame}
                Unit dari memori fisik dengan ukuran yang sama dengan page.
                \begin{itemize}
                    \item Pembagian memori fisik menjadi blok-blok tetap
                    \item Diidentifikasi dengan frame number
                    \item Setiap frame dapat menampung tepat satu page
                \end{itemize}
            \end{block}
    \end{columns}
    
    \begin{alertblock}{Penting Diingat}
        Ukuran page dan frame \textbf{harus sama} untuk memastikan proses mapping berjalan dengan benar!
    \end{alertblock}
\end{frame}

\section{Implementasi Paging}

\begin{frame}{Cara Kerja Paging}
    \framesubtitle{Proses Translasi Alamat}
    
    \begin{block}{Mekanisme Dasar}
        \begin{enumerate}
            \item CPU menghasilkan alamat logis
            \item Memory Management Unit (MMU) memecah alamat menjadi:
                \begin{itemize}
                    \item Page number (p): menunjukkan page yang diperlukan
                    \item Page offset (d): lokasi dalam page tersebut
                \end{itemize}
            \item Page number digunakan sebagai index dalam page table
            \item Page table memberikan frame number tempat page berada
            \item Frame number dan page offset digabungkan untuk membentuk alamat fisik
        \end{enumerate}
    \end{block}
    
    \begin{alertblock}{Perhatian}
        Proses translasi ini transparan bagi program yang berjalan!
    \end{alertblock}
\end{frame}

\begin{frame}{Model Paging Memori}
    \framesubtitle{Hubungan Memori Logis dan Fisik}
    
    \begin{columns}[T]
        \column{0.48\textwidth}
            \begin{block}{Memori Logis}
                \begin{itemize}
                    \item Dibagi menjadi page dengan ukuran tetap
                    \item Alamat logis: (p, d)
                    \item p = page number
                    \item d = page offset
                    \item Jumlah total page ditentukan oleh ukuran ruang alamat
                \end{itemize}
            \end{block}
            
        \column{0.48\textwidth}
                \begin{block}{Memori Fisik}
                \begin{itemize}
                    \item Dibagi menjadi frame dengan ukuran sama
                    \item Alamat fisik: $f \times$ ukuran frame + d
                    \item f = frame number
                    \item d = offset dalam frame
                    \item Frame dapat dialokasikan secara noncontiguous
                \end{itemize}
            \end{block}
    \end{columns}
    
    \begin{alertblock}{Page Table}
        Struktur data yang memetakan page ke frame, biasanya disimpan dalam memori.
        \begin{itemize}
            \item Entry ke-i berisi frame number tempat page-i disimpan
            \item Dapat menyimpan informasi tambahan (bit valid, bit proteksi, dll)
        \end{itemize}
    \end{alertblock}
\end{frame}

\begin{frame}[fragile]{Contoh Perhitungan Paging}
    \framesubtitle{Translasi Alamat Logis ke Alamat Fisik}
    
    \begin{block}{Spesifikasi Contoh}
        Memori berukuran 32 byte dengan ukuran page 4 byte:
        \begin{itemize}
            \item Memori fisik: 32 byte = 8 frame (masing-masing 4 byte)
            \item Memori logis: menggunakan page berukuran 4 byte
        \end{itemize}
    \end{block}
    
    \begin{columns}[T]
        \column{0.48\textwidth}
            \begin{alertblock}{Struktur Alamat}
                Alamat logis: 5 bit
                \begin{itemize}
                    \item 3 bit page number (0-7)
                    \item 2 bit page offset (0-3)
                \end{itemize}
            \end{alertblock}
            
        \column{0.48\textwidth}
            \begin{minted}[fontsize=\footnotesize, linenos=false]{text}
Page Table:
Page 0 → Frame 1
Page 1 → Frame 4
Page 2 → Frame 3
Page 3 → Frame 7
            \end{minted}
    \end{columns}
\end{frame}

\begin{frame}[fragile]{Contoh Perhitungan Paging (Lanjutan)}
    \framesubtitle{Translasi Alamat Logis ke Alamat Fisik}
    
    \begin{block}{Contoh Translasi}
        Alamat logis 13 (01101):
        \begin{itemize}
            \item Page number = 011 (3)
            \item Page offset = 01 (1)
            \item Page 3 disimpan di frame 7
            \item Alamat fisik = (7 $\times$ 4) + 1 = 29
        \end{itemize}
    \end{block}
\end{frame}

\begin{frame}{Ilustrasi Alokasi Frame}
    \framesubtitle{Keadaan Sebelum dan Sesudah Alokasi}
    
    \begin{columns}[T]
        \column{0.48\textwidth}
            \begin{block}{Sebelum Alokasi}
                \begin{tabular}{|c|c|}
                    \hline
                    \textbf{Frame} & \textbf{Status} \\ \hline
                    0 & Terpakai \\ \hline
                    1 & Kosong \\ \hline
                    2 & Terpakai \\ \hline
                    3 & Kosong \\ \hline
                    4 & Kosong \\ \hline
                    5 & Terpakai \\ \hline
                    6 & Terpakai \\ \hline
                    7 & Kosong \\ \hline
                \end{tabular}
                
                \scriptsize Free frames: 1, 3, 4, 7
            \end{block}
            
        \column{0.48\textwidth}
            \begin{alertblock}{Setelah Alokasi 4 Page}
                \begin{tabular}{|c|c|c|}
                    \hline
                    \textbf{Frame} & \textbf{Status} & \textbf{Page} \\ \hline
                    0 & Terpakai & - \\ \hline
                    1 & Terpakai & Page 0 \\ \hline
                    2 & Terpakai & - \\ \hline
                    3 & Terpakai & Page 2 \\ \hline
                    4 & Terpakai & Page 1 \\ \hline
                    5 & Terpakai & - \\ \hline
                    6 & Terpakai & - \\ \hline
                    7 & Terpakai & Page 3 \\ \hline
                \end{tabular}
                
                \scriptsize Semua frame sudah terpakai
            \end{alertblock}
    \end{columns}
    
    \begin{block}{Catatan}
        Alokasi frame tidak harus berurutan - paging mendukung alokasi memori noncontiguous.
    \end{block}
\end{frame}

\section{Aktivitas Kelas}

\begin{frame}{Latihan Paging}
    \framesubtitle{Implementasi Translasi Alamat (1/2)}
    
    \begin{block}{Soal}
        Diketahui sistem dengan spesifikasi:
        \begin{itemize}
            \item Ukuran memori fisik: 128 byte
            \item Ukuran page: 16 byte
            \item Page table: 
            \begin{center}
            \begin{tabular}{|c|c|}
            \hline
            \textbf{Page} & \textbf{Frame} \\
            \hline
            0 & 2 \\
            \hline
            1 & 5 \\
            \hline
            2 & 0 \\
            \hline
            3 & 6 \\
            \hline
            \end{tabular}
            \end{center}      \end{itemize}
    \end{block}
    
    \begin{alertblock}{Tugas Kelompok (5 menit)}
        \begin{enumerate}
            \item Tentukan berapa bit yang diperlukan untuk page number dan offset
        \end{enumerate}
    \end{alertblock}
\end{frame}

\begin{frame}{Latihan Paging}
    \framesubtitle{Implementasi Translasi Alamat (2/2)}
    
    \begin{alertblock}{Tugas Kelompok (lanjutan)}
        \begin{enumerate}\setcounter{enumi}{1}
            \item Terjemahkan alamat logis berikut ke alamat fisik:
            \begin{itemize}
                \item 22
                \item 53
            \end{itemize}
            \item Apakah alamat logis 73 valid? Jelaskan!
        \end{enumerate}
    \end{alertblock}
    
    \begin{block}{Petunjuk}
        Ingat rumus: Alamat fisik = (frame number $\times$ ukuran page) + offset
    \end{block}
\end{frame}

\begin{frame}{Solusi Latihan Paging}
    \framesubtitle{Analisis Kebutuhan Bit}
    
    \begin{block}{1. Kebutuhan Bit}
        \begin{itemize}
            \item Ukuran page = 16 byte → offset memerlukan 4 bit ($2^4 = 16$)
            \item Memori fisik = 128 byte → total frame = 128 / 16 = 8 frame
            \item 8 frame memerlukan 3 bit untuk representasi ($2^3 = 8$)
            \item Page number memerlukan 3 bit (sesuai dengan jumlah page)
            \item Alamat logis total: 7 bit (3 bit page number / 4 bit offset)      \end{itemize}
    \end{block}
    
    \begin{alertblock}{Struktur Alamat Logis}
        \begin{center}
            \begin{tabular}{|c|c|}
                \hline
                \textbf{Page Number} & \textbf{Offset} \\
                \hline
                3 bit & 4 bit \\
                \hline
            \end{tabular}
        \end{center}
    \end{alertblock}
\end{frame}

\begin{frame}{Solusi Latihan Paging}
    \framesubtitle{Translasi Alamat (1/2)}
    
    \begin{columns}[T]
        \column{0.48\textwidth}
            \begin{block}{Alamat logis 22}
                \begin{itemize}
                    \item Binary: 0010110
                    \item Page number: 001 (1)
                    \item Offset: 0110 (6)
                    \item Page 1 → Frame 5
                    \item Alamat fisik = (5 $\times$ 16) + 6 = 80 + 6 = 86
                \end{itemize}
            \end{block}
            
        \column{0.48\textwidth}
            \begin{block}{Alamat logis 53}
                \begin{itemize}
                    \item Binary: 0110101
                    \item Page number: 011 (3)
                    \item Offset: 0101 (5)
                    \item Page 3 → Frame 6
                    \item Alamat fisik = (6 $\times$ 16) + 5 = 96 + 5 = 101
                \end{itemize}
            \end{block}
    \end{columns}

    \begin{block}{Perhitungan Alamat Fisik}
        Alamat fisik dihitung menggunakan rumus berikut:
        \[ \text{Alamat Fisik} = (\text{Frame Number} \times \text{Ukuran Page}) + \text{Offset} \]
    \end{block}
\end{frame}

\begin{frame}{Solusi Latihan Paging}
    \framesubtitle{Translasi Alamat (2/2)}
    
    \begin{block}{Langkah Perhitungan}
        \begin{enumerate}
            \item Ekstrak page number (p) dan offset (d) dari alamat logis
            \item Temukan frame number (f) dari page table menggunakan p sebagai indeks
            \item Hitung alamat fisik dengan rumus: f × ukuran page + d
        \end{enumerate}
    \end{block}
\end{frame}

\begin{frame}{Solusi Latihan Paging}
    \framesubtitle{Validasi Alamat}
    
    \begin{block}{3. Validasi Alamat Logis 73}
        \begin{itemize}
            \item Binary: 1001001
            \item Page number: 100 (4)
            \item Offset: 1001 (9)
            \item Page table hanya mengandung Page 0-3
            \item Page 4 tidak ada dalam page table
        \end{itemize}
    \end{block}
    
    \begin{alertblock}{Kesimpulan}
        Alamat logis 73 \textbf{tidak valid}. Alasannya:
        \begin{itemize}
            \item Page number (4) melebihi rentang valid (0-3)
            \item Sistem hanya menyediakan 4 page (Page 0-3)
            \item Upaya mengakses Page 4 akan menyebabkan page fault
        \end{itemize}
    \end{alertblock}
\end{frame}

\begin{frame}{Latihan Paging Tambahan}
    \framesubtitle{Studi Kasus dengan Alamat Heksadesimal (1/2)}
    
    \begin{block}{Soal}
        Diketahui sistem dengan spesifikasi:
        \begin{itemize}
            \item Ukuran memori fisik: 4 KB (4096 byte)
            \item Ukuran page: 256 byte
            \item Page table: 
            \begin{center}
            \begin{tabular}{|c|c|}
            \hline
            \textbf{Page} & \textbf{Frame} \\
            \hline
            0 & 5 \\
            \hline
            1 & 9 \\
            \hline
            2 & 14 \\
            \hline
            3 & 3 \\
            \hline
            4 & 11 \\
            \hline
            5 & 8 \\
            \hline
            \end{tabular}
            \end{center}
        \end{itemize}
    \end{block}
\end{frame}

\begin{frame}{Latihan Paging Tambahan}
    \framesubtitle{Studi Kasus dengan Alamat Heksadesimal (2/2)}
    
    \begin{alertblock}{Tugas Kelompok (10 menit)}
        \begin{enumerate}
            \item Tentukan berapa bit yang diperlukan untuk page number dan offset
            \item Terjemahkan alamat logis berikut (dalam heksadesimal) ke alamat fisik:
            \begin{itemize}
                \item 0x2AF
                \item 0x15D
                \item 0x4C8
            \end{itemize}
            \item Apakah alamat logis 0x6A1 valid? Jelaskan!
        \end{enumerate}
    \end{alertblock}
\end{frame}

\begin{frame}{Solusi Latihan Paging Tambahan}
    \framesubtitle{Analisis Kebutuhan Bit dengan Heksadesimal}
    
    \begin{block}{1. Kebutuhan Bit}
        \begin{itemize}
            \item Ukuran page = 256 byte → offset memerlukan 8 bit ($2^8 = 256$)
            \item Memori fisik = 4096 byte → total frame = 4096 / 256 = 16 frame
            \item 16 frame memerlukan 4 bit untuk representasi ($2^4 = 16$)
            \item Page number memerlukan 3 bit (page 0-5 → maksimal 6 page → butuh 3 bit)
            \item Alamat logis total: 11 bit (3 bit page number + 8 bit offset)
        \end{itemize}
    \end{block}
    
    \begin{alertblock}{Struktur Alamat Logis}
        \begin{center}
            \begin{tabular}{|c|c|}
                \hline
                \textbf{Page Number} & \textbf{Offset} \\
                \hline
                3 bit & 8 bit \\
                \hline
            \end{tabular}
        \end{center}
    \end{alertblock}
\end{frame}

\begin{frame}{Solusi Latihan Paging Tambahan}
    \framesubtitle{Translasi Alamat Heksadesimal}
    
    \begin{columns}[T]
        \column{0.48\textwidth}
            \begin{block}{Alamat logis 0x2AF}
                \begin{itemize}
                    \item Hex: 0x2AF = 687 (desimal) = 1010101111 (biner)
                    \item Page number: 010 (2)
                    \item Offset: 10101111 (0xAF = 175)
                    \item Page 2 → Frame 14
                    \item Alamat fisik = (14 $\times$ 256) + 175 = 3584 + 175 = 3759
                \end{itemize}
            \end{block}
            
        \column{0.48\textwidth}
            \begin{block}{Alamat logis 0x15D}
                \begin{itemize}
                    \item Hex: 0x15D = 349 (desimal) = 0101011101 (biner)
                    \item Page number: 001 (1)
                    \item Offset: 01011101 (0x5D = 93)
                    \item Page 1 → Frame 9
                    \item Alamat fisik = (9 $\times$ 256) + 93 = 2304 + 93 = 2397
                \end{itemize}
            \end{block}
    \end{columns}
\end{frame}

\begin{frame}{Solusi Latihan Paging Tambahan}
    \framesubtitle{Translasi dan Validasi Alamat Hexadesimal}
    
    \begin{columns}[T]
        \column{0.48\textwidth}
            \begin{block}{Alamat logis 0x4C8}
                \begin{itemize}
                    \item Hex: 0x4C8 = 1224 (desimal) = 10011001000 (biner)
                    \item Page number: 100 (4)
                    \item Offset: 11001000 (0xC8 = 200)
                    \item Page 4 → Frame 11
                    \item Alamat fisik = (11 $\times$ 256) + 200 = 2816 + 200 = 3016
                \end{itemize}
            \end{block}
            
        \column{0.48\textwidth}
            \begin{alertblock}{Alamat logis 0x6A1 - Valid?}
                \begin{itemize}
                    \item Hex: 0x6A1 = 1697 (desimal) = 11010100001 (biner)
                    \item Page number: 110 (6)
                    \item Offset: 10100001 (0xA1 = 161)
                    \item Page table hanya sampai Page 5
                    \item Page 6 tidak terdefinisi dalam page table
                    \item \textbf{Alamat tidak valid}, akan menyebabkan page fault
                \end{itemize}
            \end{alertblock}
    \end{columns}
\end{frame}

\begin{frame}{Latihan Penerapan Paging}
    \framesubtitle{Studi Kasus Memori Berskala Besar}
    
    \begin{block}{Skenario}
        Anda sedang menganalisis sistem komputer dengan spesifikasi:
        \begin{itemize}
            \item Ruang alamat logis: 256 MB (per proses)
            \item Memori fisik: 4 GB
            \item Ukuran page: 4 KB
        \end{itemize}
    \end{block}
    
    \begin{alertblock}{Tugas Kelompok (15 menit)}
        Analisis sistem ini dengan menjawab:
        \begin{enumerate}
            \item Berapa jumlah page dan frame dalam sistem ini?
            \item Berapa bit yang diperlukan untuk page number, frame number, dan offset?
        \end{enumerate}
    \end{alertblock}
\end{frame}

\begin{frame}{Latihan Penerapan Paging (Lanjutan)}
    \framesubtitle{Studi Kasus Memori Berskala Besar}
    
    \begin{alertblock}{Tugas Kelompok (lanjutan)}
        \begin{enumerate}\setcounter{enumi}{2}
            \item Jika setiap entri page table memerlukan 4 byte, berapa total memori yang dibutuhkan untuk menyimpan page table satu proses?
            \item Jika sistem berjalan dengan 50 proses sekaligus, apakah menyimpan semua page table di memori fisik merupakan pendekatan yang efisien? Jelaskan!
        \end{enumerate}
    \end{alertblock}
    
    \begin{block}{Petunjuk Perhitungan}
        \begin{itemize}
            \item 1 KB = $2^{10}$ byte, 1 MB = $2^{20}$ byte, 1 GB = $2^{30}$ byte
            \item Jumlah page = Ukuran ruang alamat logis / Ukuran page
            \item Jumlah frame = Ukuran memori fisik / Ukuran page
        \end{itemize}
    \end{block}
\end{frame}

\begin{frame}{Solusi Latihan Penerapan Paging}
    \framesubtitle{Analisis Sistem Besar}
    
    \begin{block}{1. Jumlah Page dan Frame}
        \begin{itemize}
            \item Ukuran page = 4 KB = $2^{12}$ byte
            \item Ruang alamat logis = 256 MB = $2^{28}$ byte
            \item Jumlah page = $2^{28} \div 2^{12} = 2^{16} = 65,536$ page
            \item Memori fisik = 4 GB = $2^{32}$ byte
            \item Jumlah frame = $2^{32} \div 2^{12} = 2^{20} = 1,048,576$ frame
        \end{itemize}
    \end{block}
    
    \begin{alertblock}{2. Kebutuhan Bit}
        \begin{itemize}
            \item Page number memerlukan 16 bit (untuk $2^{16}$ page)
            \item Frame number memerlukan 20 bit (untuk $2^{20}$ frame)
            \item Offset memerlukan 12 bit (untuk 4 KB = $2^{12}$ byte)
            \item Alamat logis: 28 bit (16 bit page number + 12 bit offset)
            \item Alamat fisik: 32 bit (20 bit frame number + 12 bit offset)
        \end{itemize}
    \end{alertblock}
\end{frame}

\begin{frame}{Solusi Latihan Penerapan Paging}
    \framesubtitle{Analisis Kebutuhan Memori}
    \begin{columns}[T]
        \column{0.48\textwidth}
            \begin{block}{3. Ukuran Page Table}
                \begin{itemize}
                    \item Jumlah entri page table = 65,536 entri
                    \item Ukuran per entri = 4 byte
                    \item Total ukuran page table = $65,536 \times 4 = 262,144$ byte = 256 KB
                \end{itemize}
            \end{block}
            
        \column{0.48\textwidth}
            \begin{alertblock}{4. Efisiensi Penyimpanan Page Table}
                \begin{itemize}
                    \item 50 proses membutuhkan $50 \times 256$ KB = 12,800 KB = 12.5 MB untuk page table
                    \item 12.5 MB dari 4 GB (0.3\% dari memori fisik)
                    \item Secara persentase, ini kecil, tapi masih signifikan
                    \item \textbf{Tidak efisien} karena:
                    \begin{itemize}
                        \item Tidak semua page digunakan oleh proses
                        \item Overhead terus meningkat dengan jumlah proses
                    \end{itemize}
                \end{itemize}
            \end{alertblock}
    \end{columns}
\end{frame}

\begin{frame}{Solusi Latihan Penerapan Paging}
    \framesubtitle{Alternatif Desain dan Implikasi}
    
    \begin{block}{Alternatif Solusi}
        \begin{itemize}
            \item Implementasi page table bertingkat (multilevel page table)
            \item Hanya menyimpan bagian page table yang aktif digunakan
            \item Menggunakan Translation Lookaside Buffer (TLB) untuk mempercepat akses
            \item Implementasi page table terbalik (inverted page table)
        \end{itemize}
    \end{block}
    
    \begin{alertblock}{Implikasi Penting}
        \begin{itemize}
            \item Meski persentase memori untuk page table kecil, penggunaan teknik yang efisien tetap penting
            \item Dengan ruang alamat yang lebih besar (64-bit), masalah ini semakin signifikan
            \item Solusi seperti multilevel paging dan TLB menjadi sangat penting dalam sistem modern
            \item Trade-off: kompleksitas implementasi vs efisiensi memori
        \end{itemize}
    \end{alertblock}
\end{frame}

\section{Kesimpulan}

\begin{frame}{Rangkuman Paging}
    \framesubtitle{Poin-poin Penting}
    
    \begin{columns}[T]
        \column{0.48\textwidth}
            \begin{block}{Konsep Utama}
                \begin{itemize}
                    \item Paging membagi memori logis menjadi page berukuran tetap
                    \item Memori fisik dibagi menjadi frame dengan ukuran sama
                    \item Page table menangani translasi alamat
                    \item Mendukung alokasi memori noncontiguous
                \end{itemize}
            \end{block}
            
        \column{0.48\textwidth}
            \begin{alertblock}{Tantangan Implementasi}
                \begin{itemize}
                    \item Ukuran page table bisa menjadi besar
                    \item Overhead untuk translasi alamat
                    \item Memerlukan hardware support (MMU)
                    \item Perlu mekanisme proteksi dan sharing
                \end{itemize}
            \end{alertblock}
    \end{columns}
    
    \begin{block}{Pengembangan}
        Teknik paging adalah fundamental untuk implementasi memori virtual, yang akan kita pelajari selanjutnya.
    \end{block}
\end{frame}

\backmatter % generate the final slide
\end{document}
